\subsection{Problem Setting}

We study binary still-image classification (label $0$: real; label $1$: AI-generated) under compact-model constraints.
Video deepfakes, legal-forensics certification, and large-scale foundation-model detectors are out of scope.
The intended use case is offline or server-side screening where Params/FLOPs/latency and honest OOD boundaries matter alongside ID AUC.

\subsection{Datasets and Splits}

\paragraph{Training and ID evaluation}
use DiffusionForensics (DF) images from two content domains present in training: ADM bedroom and Stable Diffusion~v2 CelebA-HQ, each with real and fake folders.
Frozen JSONL manifests (seed \textbf{42}) contain train/val/test/ood splits of sizes \textbf{3200/400/400/5600}.
No test or OOD image is used for training, validation AUC selection, or threshold fitting.

\paragraph{UniversalFakeDetect (UFD)}
provides the primary cross-generator suite~\cite{ojha2023ufd}.
For each subset we sample up to 400 real and 400 fake images (seed 42).
Subsets in the frozen package include \texttt{dalle}, \texttt{glide\_100\_10}, \texttt{glide\_100\_27}, \texttt{glide\_50\_27}, \texttt{guided}, \texttt{ldm\_100}, \texttt{ldm\_200}, and \texttt{ldm\_200\_cfg}.

\paragraph{Exploratory FLUX set (appendix only).}
Independently of training, we generated 100 images with \texttt{black-forest-labs/FLUX.1-schnell}~\cite{flux2024} and paired them with 100 DF real images (seed 42).
Because real/fake subsets were not content- and acquisition-matched, results appear only in Appendix~\ref{app:flux} and are never used for model selection or architecture claims.

\subsection{Preprocessing}

All pipelines resize images to $224\times224$, convert to RGB, and apply ImageNet mean/std normalization.
Training adds random horizontal flip; evaluation is deterministic.
JPEG experiments recompress decoded RGB images to JPEG quality $Q\in\{95,85,70\}$ in memory before the same resize/normalize stack; the main robustness table freezes \textbf{Q70} versus clean.

\subsection{Models}

\begin{itemize}
\item EfficientNet-B0, MobileNetV3-Small, ShuffleNetV2-x0.5: \texttt{torchvision} classifiers with \texttt{weights=None}~\cite{tan2019efficientnet,howard2019mobilenetv3,ma2018shufflenet}.
\item LiteFreqNet~v2: spatial stem + FFT-magnitude branch with gated residual fusion $\mathrm{LN}(s+\sigma(\alpha)f)$ and mid-band prior; early-concat and no-frequency variants used as ablations.
\item Two study-specific SSM classifiers, detailed in Section~\ref{sec:litessm}.
\end{itemize}
No ImageNet pretrained weights are used for any Panel~A model.

\subsection{Study-Specific SSM Classifiers (LiteSSM-A / LiteSSM-B)}
\label{sec:litessm}

\textbf{LiteSSM-A} denotes the compact baseline SSM classifier, whereas \textbf{LiteSSM-B} extends it with a BiViM-inspired bidirectional SelectiveSSM design.
Both models are study-specific pure-PyTorch implementations rather than official reproductions of an existing architecture.
Frozen experiment artifacts retain the registry keys \texttt{mobilemamba\_lite} (LiteSSM-A) and \texttt{mambapsa\_cls} (LiteSSM-B) so that checkpoints and hashes remain unchanged; the manuscript uses only the LiteSSM names.
Optional \texttt{*\_freq} DualBranch wrappers attach the gated frequency branch for ablation only and are not main-table models.

\paragraph{Shared pipeline.}
Inputs are $B\times3\times224\times224$.
A four-stage stride-2 convolutional stem (Conv $3\times3$, BatchNorm, GELU) yields a $14\times14\times192$ map, which is flattened to $L=196$ tokens of width $C=192$, summed with a learnable positional embedding, processed by depth~$=4$ residual blocks, LayerNormed, mean-pooled over tokens, and classified by a linear head to two logits.

\paragraph{Selective scan without vendor kernels.}
Both models use a pure-PyTorch \texttt{SelectiveSSM}: input projection, depthwise Conv1d ($k=3$), input-dependent $\Delta,B,C$, and an explicit sequential recurrence over $L$ (no CUDA/\texttt{mamba\_ssm} kernel).
This choice keeps Params/FLOPs/latency comparisons portable but can inflate absolute latency relative to fused official kernels; reported batch-1 times are therefore \emph{model-only measurements under the pure-PyTorch stack}.

\paragraph{LiteSSM-A block (MRFFILite).}
After LayerNorm, a multi-receptive local branch (depthwise Conv1d kernels $\{3,5,7\}$ summed) is concatenated with a SelectiveSSM global branch ($d_{\mathrm{state}}=8$, expand~$=1$), projected by a linear layer, and added residually.

\paragraph{LiteSSM-B block (BiViM).}
After LayerNorm, forward and time-reversed SelectiveSSM branches ($d_{\mathrm{state}}=16$, expand~$=1$) are concatenated, linearly mixed, and added residually.

\begin{table}[H]
\caption{Architecture configuration of the study-specific SSM models.
Both use input $224\times224$, stem downsampling $\times16$, token length $L=196$, width $C=192$, and depth~$=4$ residual blocks.\label{tab:ssm-arch}}
\begin{tabularx}{\textwidth}{l R l l}
\toprule
\textbf{Stage} & \textbf{Out size} & \textbf{LiteSSM-A} & \textbf{LiteSSM-B} \\
\midrule
Stem-1 & $112\times112$ & Conv/s2, BN, GELU (64) & Conv/s2, BN, GELU (96) \\
Stem-2 & $56\times56$ & Conv/s2, BN, GELU (128) & Conv/s2, BN, GELU (96) \\
Stem-3 & $28\times28$ & Conv/s2, BN, GELU (192) & Conv/s2, BN, GELU (192) \\
Stem-4 & $14\times14$ & Conv/s2, BN, GELU (192) & Conv/s2, BN, GELU (192) \\
Blocks & $L{=}196$ & $4\times$ MRFFILite (SSM $d_s{=}8$) & $4\times$ BiViM (SSM $d_s{=}16$) \\
Head & $1$ & LN, mean pool, Linear$\to$2 & LN, mean pool, Linear$\to$2 \\
\midrule
Params / FLOPs & -- & 1.74\,M / 0.705\,G & 2.48\,M / 0.853\,G \\
\bottomrule
\end{tabularx}
\end{table}

\begin{figure}[H]
\centering
\includegraphics[width=0.98\textwidth]{figures/fig_ssm_architecture.png}
\caption{Overview of LiteSSM-A and LiteSSM-B.
Both share a $\times16$ stem and $L{=}196$ token interface; blocks differ (multi-receptive+SSM vs.\ bidirectional SSM).
No official CUDA selective-scan kernel is used.\label{fig:ssm-arch}}
\end{figure}

\paragraph{Efficiency accounting.}
Trainable parameter counts sum \texttt{requires\_grad} tensors; FLOPs use \texttt{thop} on a $1\times3\times224\times224$ tensor.
All reported efficiency numbers come from one locked measurement session on an RTX~4090D (PyTorch~2.3, CUDA~12.1, FP32): batch-1 latency with $\ge$50 warmup and 500 timed iterations under \texttt{torch.cuda.synchronize()}, and batch-32 throughput in the same session.
Latency/throughput are \textbf{model-only} on synthetic tensors and exclude disk I/O and torchvision preprocessing; FFT modules are included when present in the evaluated network.

\subsection{Training Recipe (Locked)}

Every main-table model uses 15 epochs, batch size 64, AdamW ($\mathrm{lr}=10^{-4}$, weight decay $10^{-4}$), CosineAnnealingLR, cross-entropy loss, input size 224, and seed 42 on one NVIDIA GeForce RTX~4090D (PyTorch~2.3).
After each epoch we keep the checkpoint with the best validation AUC.
A decision threshold is then fit on validation probabilities by maximizing Youden's $J=\mathrm{TPR}-\mathrm{FPR}$.
Primary ranking metrics in this paper are AUC-based.

\subsection{Metrics and Statistical Protocol}
\label{sec:metrics}

\textbf{UFD Macro AUC} (primary OOD) equals the equal-weight mean of per-generator AUCs on UFD subsets.
\textbf{Worst-generator AUC} is the minimum subset AUC.
\textbf{OOD Pooled AUC} (supplementary) is a single AUC on all rows of \texttt{test\_ood.jsonl} and is \emph{not} interchangeable with UFD Macro.
Domain AUCs partition the ID test set by \texttt{celebahq} vs.\ \texttt{bedroom} source strings.
JPEG $\Delta$AUC is $\mathrm{AUC}_{Q70}-\mathrm{AUC}_{\mathrm{clean}}$ on the ID test set.

Efficiency metrics are taken exclusively from the locked session summarized in \texttt{freeze/frozen\_numbers.json}: Params, FLOPs, batch-1 p50/p95 latency, and batch-32 throughput.
Preliminary throughput figures produced outside that protocol are superseded and are not reported in the manuscript.

Unless stated otherwise, percentile 95\% confidence intervals use $B=1000$ resamples with seed~42 on frozen per-sample scores.
For UFD Macro, each bootstrap redraws within every generator subset then re-averages.
Paired $\Delta$AUC CIs redraw the same indices for both models.

\subsection{JPEG Evaluation Protocol}

Clean and JPEG evaluations share the ID test manifest.
For quality $Q$, each image is loaded as RGB, encoded to JPEG at quality $Q$, decoded, then passed through the evaluation transform.
We report Q70 as the primary robustness operating point in Table~\ref{tab:jpeg}.

\subsection{Exploratory FLUX Generation Protocol}
\label{sec:flux-protocol}

Protocol details and the unmatched-subset caveat are given in Appendix~\ref{app:flux}; the main text does not use FLUX scores for architecture claims.

\subsection{External Reference Detectors (Inference Only)}
\label{sec:external-refs}

To contextualize Panel~A without breaking the from-scratch fairness claim, we evaluate two publicly released detectors on the same frozen manifests using their native preprocessing and pretrained weights:
\begin{itemize}
\item UnivFD (CLIP ViT-L/14 linear probe)~\cite{ojha2023ufd};
\item NPR (neighboring-pixel-relationship CNN)~\cite{tan2024npr}.
\end{itemize}
DIRE~\cite{wang2023dire} is discussed in Related Work only and is not re-run.
Results form Panel~B and are captioned as references, not training-matched competitors.

\subsection{Implementation Artifacts}

Training/evaluation code, freeze/bootstrap scripts, prediction shards, table exports, and SHA256 digests of manifests/checkpoints/predictions are retained so statistics can be recomputed without retraining.
A public repository and archival DOI are planned before submission; Data Availability will list original dataset providers, license limits on redistribution, manifest hashes, and code/weight locations.
